\documentclass[12pt, letterpaper]{article}

%-------Packages---------
\usepackage{newpxtext,newpxmath} % Palatino-like text & math (modern)
\usepackage{bboldx}
\usepackage{mathtools}
\usepackage{tikz}
\usetikzlibrary{calc,intersections}
\usepackage{tikz-cd}
\usepackage[all,arc]{xy}
\usepackage{graphicx}

\usepackage{enumerate}
\usepackage[dvipsnames]{xcolor}
\usepackage{float}
\usepackage[left=1in,top=1in,right=1in,bottom=1in]{geometry}
\usepackage{fancyhdr}
\usepackage{multicol}
\usepackage{setspace}

\usepackage{dsfont}
\usepackage{mathrsfs}

\usepackage{tcolorbox}
\tcbuselibrary{theorems,skins,breakable}

\usepackage{xparse}
\usepackage{import}
\usepackage[utf8]{inputenc}

\usepackage{hyperref}
\usepackage{cleveref}

\usepackage{xcolor, ifthen, tcolorbox}
\tcbuselibrary{theorems,skins,breakable}
% Theorem System
% The following environments are provided:
%   New Problem:        \begin{problem} ...
%   New Question Header:   \qh (then followed by subproblems)
%   Subproblem:     \begin{subproblem} ...
%   Problem Proof:  \begin{problemproof} ...
%   Proof:          \\begin{pf}{color} ...
%   Remark:         \rmk (sentence), \rmkb (block)

% Define custom colors
\definecolor{thmscol}{HTML}{F4565B} %For Problems
\definecolor{lemscol}{HTML}{FE844C} %For Lemmes
\definecolor{prpscol}{HTML}{00A7EF} %For proposition
\definecolor{corscol}{HTML}{23DAFF} %For corrolaries
\definecolor{rmkscol}{HTML}{6DEEC5} %For remarks
\definecolor{defscol}{HTML}{18C991} %For definitions
\definecolor{exmscol}{HTML}{7DB800} %For examples
\definecolor{clmscol}{HTML}{165c58} %For claims

%Change between dark(black)/light(white) mode
\newcommand{\colormode}{white}
\ifthenelse{\equal{\colormode}{white}}
      {\newcommand{\TextColor}{black}}
      {\newcommand{\TextColor}{white}}
\newcommand{\pbcolormain}{thmscol!80!\colormode}
\newcommand{\pbcolorsecond}{thmscol!38!\colormode}


% ============================
% Problem Counter
% ============================
\newcounter{subproblemcounter}
\newcounter{problemcounter}

\newcommand{\q}{
    \setcounter{subproblemcounter}{0}%
    \stepcounter{problemcounter} % Increment problem counter
    \color{\TextColor}Problem \arabic{problemcounter}: % Display the problem counter
}

\newcommand{\stepsub}{
    \stepcounter{subproblemcounter}%
    \color{\TextColor}(\alph{subproblemcounter})
}
% ============================
% Problem
% ============================
\newtcolorbox{pb}[1]{
  enhanced,
  frame hidden,
  titlerule=0mm,
  toptitle=1mm,
  bottomtitle=1mm,
  fonttitle=\bfseries\large,
  coltitle=black,
  colbacktitle=\pbcolormain,
  colback=\pbcolorsecond,
  title={\q #1},
  coltext=\TextColor
}

\newtcolorbox{spb}[1]{
  enhanced,
  frame hidden,
  titlerule=0mm,
  toptitle=1mm,
  bottomtitle=1mm,
  fonttitle=\bfseries\large,
  coltitle=black,
  colbacktitle=\pbcolormain,
  colback=\pbcolorsecond,
  title={\stepsub #1},
  coltext=\TextColor,
  attach boxed title to top left={xshift=3mm,yshift=-2mm},
  boxed title style={
    colback=\pbcolormain,
    colframe=\pbcolormain,
  }
}

\newtcolorbox{pbh}[1]{
    enhanced,
    frame hidden,
    titlerule=0mm,
    toptitle=1mm,
    bottomtitle=1mm,
    fonttitle=\bfseries\large,
    coltitle=black,
    colbacktitle=\pbcolormain,
    colback=\pbcolorsecond,
    title={\q #1},
    boxrule=0pt,
    interior hidden,
    attach boxed title to top left={xshift=3mm,yshift=-2mm},
    boxed title style={
        colback=\pbcolormain,
        colframe=\pbcolormain,
    }
}

\newenvironment{problem}[1]{%
  \newpage
  \begin{pb}{#1}
    }{
  \end{pb}
}

\newenvironment{subproblem}[1]{%
  \begin{spb}{#1}
    }{
  \end{spb}
}

\newenvironment{problemproof}{%
    \begin{pf}{\pbcolormain}
        }{
    \end{pf}
}

\newcommand{\qh}[1][]{
    \newpage
    \begin{pbh}{#1}\end{pbh} % Display the problem counter
}

% ============================
% Claim
% ============================

\newtcbtheorem[number within=problemcounter]{claim}{Claim}
{
    enhanced,
    frame hidden,
    titlerule=0mm,
    toptitle=1mm,
    bottomtitle=1mm,
    fonttitle=\bfseries\large,
    coltitle=black,
    colbacktitle=clmscol!50!\colormode,
    colback=clmscol!20!\colormode,
}{clm}

\NewDocumentCommand{\clm}{m+m}{
    \begin{claim*}{#1}{}
        #2
    \end{claim*}
}

\newenvironment{clmpf}{
	{\noindent{\it \textbf{Proof.}}
	\setlength{\parindent}{0pt}
	}
	\tcolorbox[blanker,breakable,left=5mm,parbox=false,
    before upper={\parindent15pt},
    after skip=10pt,
	borderline west={1mm}{0pt}{clmscol!50!\colormode}]
}{
    \textcolor{clmscol!50!\colormode}{\hbox{}\nobreak\hfill$\blacksquare$}
    \endtcolorbox
}

\NewDocumentCommand{\clmp}{m+m+m}{
    \begin{claim*}{#1}{}
        #2
    \end{claim*}

    \begin{clmpf}
        #3
    \end{clmpf}
}


% ============================
% Proof
% ============================

\newenvironment{pf}[1]{%
  \def\pfcolor{#1}% Store the color in a macro
  \noindent{\it \textbf{Proof.}}%
  \tcolorbox[blanker,breakable,left=5mm,parbox=false,%
    before upper={\parindent15pt},%
    after skip=10pt,%
    borderline west={1mm}{0pt}{\pfcolor},%
    coltext=\TextColor
  ]%
  \setlength{\parindent}{0pt}%
}{%
  \textcolor{\pfcolor}{\hbox{}\nobreak\hfill$\blacksquare$}%
  \endtcolorbox%
}

\newtcolorbox{solution}{
  blanker,
  breakable,
  left=5mm,
  parbox=false,%
  before upper={\parindent15pt},%
  after skip=10pt,%
  borderline west={1mm}{0pt}{\pbcolormain},%
  coltext=\TextColor
}


% ============================
% Remark
% ============================


\NewDocumentCommand{\rmk}{+m}{
    {\it \color{rmkscol!100!white}#1}
}

\newenvironment{remark}{
    \par
    \vspace{5pt}
    \begin{minipage}{\textwidth}
        {\par\noindent{\textbf{Remark.}}}
        \tcolorbox[blanker,breakable,left=5mm,
        before skip=10pt,after skip=10pt,
        borderline west={1mm}{0pt}{rmkscol!80!\colormode},
        coltext=\TextColor
        ]
}{
        \endtcolorbox
    \end{minipage}
    \vspace{5pt}
}

\NewDocumentCommand{\rmkb}{+m}{
    \begin{remark}
        #1
    \end{remark}
}


% Define a new command for problems
\renewcommand{\headrule}{\color{\TextColor}\hrulefill}
\newcommand{\C}{\mathbb{C}}
\newcommand{\R}{\mathbb{R}}
\newcommand{\Q}{\mathbb{Q}}
\newcommand{\Z}{\mathbb{Z}}
\newcommand{\N}{\mathbb{N}}
\newcommand{\p}{\mathfrak{p}}
\newcommand{\F}{\mathbb{F}}
\newcommand{\contr}{\Rightarrow \Leftarrow}
\newcommand{\s}{\(\,\)}
\renewcommand{\t}[1]{\text{#1}}
\newcommand{\Spec}{\mathrm{Spec}}
\newcommand{\Ass}{\mathrm{Ass}}
\newcommand{\Supp}{\mathrm{Supp}}
\newcommand{\Ann}{\mathrm{Ann}}
\newcommand{\mf}[1]{\mathfrak{#1}}
\newcommand{\codim}{\mathrm{codim}}

% Meta data
\setstretch{1.2}

\begin{document}
\color{\TextColor}
\pagecolor{\colormode}
\setlength{\parindent}{0pt}
\pagestyle{fancy}
\fancyhf{}
\fancyhead[LOE]{\color{\TextColor}\slshape{\textbf{Vincent Ferrara}}}
\fancyhead[ROE]{\color{\TextColor}HW11 --- \thepage}
\fancyhead[C]{\color{\TextColor}Math 614}

\begin{problem}{}
    Prove that if $R$ is a valuation ring, then $R$ is normal.
\end{problem}
\begin{problemproof}
    Let $R$ be a valuation ring.

    Let $x \in K=\mathrm{Frac}(R)$ be integral over $R$. This means $x$ is a root of a monic polynomial.

    $x^n+a_{n-1}x^{n-1}+\dots+a_0=0$ for $a_i \in R$.

    Assume for contradiction that $x \notin R$. This implies that $x^{-1} \in R$ by definition of a valuation ring.

    This implies that $x^n=-(a_{n-1}x^{n-1}+\dots+a_0) \implies x=(-a_{n-1}+a_{n-1}x^{-1}+\dots+x_0x^{-(n-1)})$.

    This implies that $x \in R$.

    Thus, $R$ is normal.
\end{problemproof}

\begin{problem}{}
    Let $(R, \mf{m})$ be a noetherian local ring with residue field $k = R \diagup \mf{m}$.
\end{problem}
\begin{subproblem}{}
    Prove that $\mf{m}$ is principal if and only if $\mf{m} \diagup \mf{m}^2$ has dimension $\leq 1$ as a $k$-vector space.
\end{subproblem}
\begin{problemproof}
    Assume $\mf{m}=(x)$ for some $x \in \mf{m}$.

    If $x = 0$ then $\mf{m}=0 \implies \mf{m} \diagup \mf{m^2} = 0 \implies$ it has dimension 0.

    If $x \neq 0$ then this implies that $\bar{x}$ which is the image of $x$ onto $\mf{m} \diagup \mf{m}^2$ generates this ideal.

    Let $\bar{y} \in \mf{m} \diagup \mf{m}^2 \implies \pi(y)=\bar{y}$. Thus, $y=rx$ for some $r \in R$. Thus, $\pi(y)=\pi(r)\pi(x)=\bar{r}\bar{x}$.

    Therefore, $\bar{y}=\bar{r}\bar{x}$. Thus, $\mf{m} \diagup \mf{m}^2 = (\bar{x})$. Thus, this has dimension 1.

    Assume $\mf{m} \diagup \mf{m}^2$ has dimension $\leq 1$ has a $k$-vector space.

    If $\dim (\mf{m} \diagup \mf{m}^2) = 0$

    Then $\mf{m} \diagup \mf{m}^2 = 0 \implies \mf{m} = \mf{m}^2$.

    Since $\mf{m}$ is finitely generated, by Nakayama's lema we have that $\mf{m}=0=(0)$.

    If $\dim (\mf{m} \diagup \mf{m}^2) = 1$

    Then there exists some $x \in \mf{m}$ s.t. $\bar{x}$ generates $\mf{m} \diagup \mf{m}^2=(\bar{x})$.

    Consider $(x)$

    By construction, the image of $(x)$ in $\mf{m} \diagup \mf{m}^2$ generate $\mf{m} \diagup \mf{m}^2$.

    Thus, $\mf{m}=(x)+\mf{m}^2$.

    By Nakayama's lemma since $\mf{m}$ is finitely generated since $R$ is noetherian we get that $\mf{m}=(x)$.

    Thus, $\mf{m}$ is principal in both cases.
\end{problemproof}

\begin{subproblem}{}
    Prove that $R$ is a DVR if and only if $R$ is a domain and $\mf{m} \diagup \mf{m}^2$ has dimension 1 as a
$k$-vector space.
\end{subproblem}
\begin{problemproof}
    Assume $R$ is a DVR. This implies $R$ is a PID from shown in class. Thus, by (a) this implies that $\dim \mf{m} \diagup \mf{m}^2 \leq 1$, but since $\mf{m} \neq 0$ since $R$ is not a field this implies that $\dim \mf{m} \diagup \mf{m}^2=1$.

    Also $R$ is a domain since it is a PID.

    Assume $R$ is a domain and $\mf{m} \diagup \mf{m}^2$ has dimension 1 as a $k$-vector space.

    By (a) this implies that $\mf{m}=(x)$ for some $x \in R$. But since $\mf{m} \diagup \mf{m}^2$ has dimension one this implies $\mf{m} \neq 0$.

    Shown in class this implies that $R$ is a DVR.
\end{problemproof}

\begin{subproblem}{}
    Give a counter example where $\mf{m} \diagup \mf{m}^2$ has dimension 1 as a $k$-vector space, but $R$ is not a DVR.
\end{subproblem}
\begin{problemproof}
    Consider $R= k[x] \diagup (x^2)$ for field $k$.

    $R$ is noetherian and local with maximal $\mf{m}=(\bar{x})$.

    Consider $(\bar{x}) \diagup (\bar{x}^2) = (\bar{x}) \diagup (0) \cong (\bar{x})$. Thus, $\dim(\mf{m} \diagup \mf{m}^2)=1$.

    But $R$ is not a DVR since it is not a domain.
\end{problemproof}

\begin{problem}{}
    Let $k$ be a field and let $f \in k[x, y]$ be an irreducible polynomial of the form $f = \ell + g$ where
$\ell = ax + by$ is the linear part and $g \in (x, y)^2$. Let $R = k[x, y]/(f )$, let $m = (\bar{x}, \bar{y}) \subset R$ be
the maximal ideal generated by the images of $x$ and $y$ in R. Prove that the localization $R_\mf{m}$
is a DVR if and only if $\ell \neq 0$.
\end{problem}
\begin{problemproof}
    Since $k[x,y]$ is a finite type $k$ algebra this implies that $\dim k[x,y] = \dim R + \codim (f) \implies \dim R = 1$. This is because $\dim k[x,y]=2$ and by the principal ideal theorem in a domain we have $\codim (f) = 1$.
\end{problemproof}

\begin{problem}{}
    Prove that if $R$ is a DVR with maximal ideal $\mf{m}$, then the $\mf{m}$-adic completion $\widehat{R}$ of $R$ is a
DVR.
\end{problem}
\begin{problemproof}
    Assume $R$ is a DVR with maximal ideal $\mf{m}$. Let $\mf{m}=(\pi)$.
\end{problemproof}

\begin{problem}{}
    In this problem, we show that normal noetherian rings decompose as products of normal
noetherian domains.
\end{problem}
\begin{subproblem}{}
    Let $R$ be a noetherian ring such that $R_\p$ is a domain for every $\p \in \Spec(R)$. Prove that
there exist domains $R_1, \dots , R_n$ such that there is an isomorphism $R \cong R_1 \times \dots \times R_n$.
\end{subproblem}
\begin{problemproof}
    Let $x \in \widehat{R}$. Thus, $x=(x_i)$ s.t. $x_i = x_{i+1} \mod \mf{m}^i$ for all $i$.

    Show that $x= \sum_{i=0}^\infty a_i \pi^i$.

    Base Case: $n=1$.

    $x_1 \in R \diagup (\pi)$. Choose $a_0 \in R$ whoes image is $x_1$.

    Inductive Step: Suppose for some $n \geq 1$ we have chosen $a_0,\dots,a_{n-1} \in R $ s.t. 
    
    $x_n \equiv \sum_{i=0}^{n-1} a_i \pi^i \pmod{\pi^n}$.

    Consider $x_{n+1} \in R \diagup (\pi^{n+1})$. It's image in $R \diagup (\pi^n)$ is $x_n$, and the image of $\sum_{i=0}^{n-1} a_i \pi^i$ in $R \diagup (\pi^n)$ is also $x_n$.

    Thus,
    \[x_{n+1} - \sum_{i=0}^{n-1} a_i \pi^i \in (\pi^n) \diagup (\pi^{n+1})\]
    
\end{problemproof}
\end{document}